\documentclass[11pt]{article}
\usepackage[a4paper,margin=27mm]{geometry}
\usepackage{amsmath,amssymb,amsthm,mathtools,bm,booktabs,longtable,microtype}
\usepackage[T1]{fontenc}
\usepackage{lmodern}
\usepackage[hidelinks]{hyperref}
\usepackage{enumitem}
\setlist{nosep}
\newtheorem{theorem}{Theorem}[section]
\newtheorem{lemma}[theorem]{Lemma}
\newtheorem{proposition}[theorem]{Proposition}
\newtheorem{corollary}[theorem]{Corollary}
\theoremstyle{remark}\newtheorem{remark}[theorem]{Remark}
\newcommand{\R}{\mathbb R}\newcommand{\Ssym}{\mathbb S}\newcommand{\K}{\mathcal K}
\newcommand{\ip}[2]{\langle #1,#2\rangle}\newcommand{\norm}[1]{\lVert #1\rVert}
\newcommand{\tr}{\operatorname{tr}}\newcommand{\dist}{\operatorname{dist}}
\newcommand{\Clust}{\operatorname{Clust}}\newcommand{\diag}{\operatorname{diag}}
\newcommand{\ran}{\operatorname{ran}}\newcommand{\Fix}{\operatorname{Fix}}
\newcommand{\eps}{\varepsilon}\newcommand{\dd}{\Delta}
\DeclareMathOperator*{\argmin}{argmin}
\numberwithin{equation}{section}
\title{Static excess and spectral transport at infinity\newline in conic Douglas--Rachford splitting}
\author{Research Track I: working manuscript}
\date{15 November 2026}
\begin{document}
\maketitle
\begin{abstract}
For an affine space and a closed convex semialgebraic cone disjoint at distance zero, plain Douglas--Rachford iterates have a common leading state $Vk^p$. That state need not determine their leading velocity. We identify a nonnegative static excess of the cone shadow whose macroscopic averages equal half the transverse kinetic energy. Vanishing excess is equivalent to the exact last-step law, whereas a positive limiting excess determines both the excess speed and logarithmically divergent spherical length. For products of positive semidefinite cones, an interior leading state supplies a strictly contracting normal spectral dynamics and forces the exact square-root velocity. In contrast, a rational order-seven affine family has a circular scaled last-step limit set, an exact nonradial speed coefficient, and infinite spherical length. Its proof uses an exact matrix-cut energy, a weighted spectral chart, stable fast variables, a periodic filter, and a quartic amplitude correction. The resonance persists on a rational parameter interval and on relatively open sets of starts in an invariant block space. A frozen switching map explains why a previously simulated parameter value is not covered by the separated spectral argument.
\end{abstract}
\paragraph{Status.} This is a research draft, not a submission-cleared manuscript. The primary-literature and journal-comparison audits are incomplete. The original small rational start at parameter $\gamma=1$ is not claimed to be resolved. The counterexample proved here uses $3/2\leq\gamma\leq5/2$ and an explicitly specified rational starting family. All computations used in the proofs are exact; finite numerical trajectories are not proof premises.
\tableofcontents

\section{Problem and main conclusions}
Let $E$ be a finite-dimensional Euclidean space, let $\K\subset E$ be a closed convex semialgebraic cone, and let $U=a+L$ be a nonempty affine space, where $a\perp L$. Assume
\begin{equation}\label{eq:weak}
 U\cap\K=\varnothing,\qquad \dist(U,\K)=0.
\end{equation}
Write $P=P_L$, $Q=I-P$, and consider precisely the affine-first, unanchored iteration
\begin{equation}\label{eq:drs}
 Tz=z+P_\K(2P_Uz-z)-P_Uz,\qquad z_{k+1}=Tz_k.
\end{equation}
Every projection uses the given Euclidean metric. The positive dual cone is $\K^*=\{s:\ip{s}{v}\geq0\ (v\in\K)\}$.

The map $T$ is semialgebraic and firmly nonexpansive. It has no fixed point and its infimal displacement is zero. Section~\ref{sec:state} proves the existence of an initial-point-independent nonzero $V$ and rational $p\in(0,1)$ such that
\begin{equation}\label{eq:state}
 z_k=Vk^p+o(k^p).
\end{equation}
The question is whether $k^{1-p}(z_{k+1}-z_k)\to pV$. It is essential not to differentiate the remainder in \eqref{eq:state}.

Put $z_k=a+x_k+y_k$, with $x_k=Pz_k$ and $y_k=Qz_k-a$, and set
\begin{equation}\label{eq:shadows}
 d_k=z_{k+1}-z_k,\quad v_k=P_\K(a+x_k-y_k),\quad
 s_k=P_{\K^*}(-a-x_k+y_k).
\end{equation}
Define $H=PV$, $G=QV$ and the static distance profile
\begin{equation}\label{eq:profile}
 D(r)=\min\{\norm{Qv-a}:v\in\K,\ \norm{Pv}\leq r\}.
\end{equation}
For the actual cone shadows define
\begin{equation}\label{eq:excess}
 r_k=\norm{Pv_k},\quad \delta_k=\norm{Qv_k-a},\quad
 e_k=\delta_k^2-D(r_k)^2\geq0.
\end{equation}
For $\lambda>1$, let $m=\lfloor\lambda n\rfloor$ and
\[
 I_{n,m}=\sum_{j=n}^{m-1}j^{2p-2},\qquad
 X_n(\lambda)=I_{n,m}^{-1}\sum_{j=n}^{m-1}e_j,
 \qquad A_0=p^2\norm V^2.
\]

\begin{theorem}[Static excess and transverse kinetic energy]\label{thm:excess}
For every cone and affine space satisfying \eqref{eq:weak}, the bound $d_k=O(k^{p-1})$ is equivalent to boundedness of $X_n(\lambda)$ for one, equivalently every, fixed $\lambda>1$. Under these equivalent conditions,
\begin{align}
 k^{1-p}\ip{d_k}{V}&\longrightarrow p\norm V^2,\label{eq:radial}\\
 \sum_{j=n}^{m-1}\norm{d_j}^2
 &=A_0I_{n,m}+2\sum_{j=n}^{m-1}e_j+o(I_{n,m}),\label{eq:energy}\\
 \frac{\sum_{j=n}^{m-1}j^{2p-2}\norm{j^{1-p}d_j-pV}^2}{I_{n,m}}
 &=2X_n(\lambda)+o(1).\label{eq:transverse}
\end{align}
The following are equivalent, without separately assuming the bound on $d_k$:
\begin{equation}\label{eq:equiv}
 k^{1-p}d_k\to pV;\qquad
 X_n(\lambda)\to0\text{ for one fixed }\lambda>1;\qquad
 \frac{\delta_k}{D(r_k)}\to1.
\end{equation}
For $e\geq0$, convergence $X_n(\lambda)\to e$ for every fixed $\lambda>1$ is equivalent to
\begin{equation}\label{eq:normlimit}
 k^{2-2p}\norm{d_k}^2\longrightarrow A_0+2e.
\end{equation}
If these hold and $\nu_k=z_k/\norm{z_k}$, then
\begin{align}
 \frac{\sum_{k=K}^{n-1}\norm{\nu_{k+1}-\nu_k}}{\log n}
 &\longrightarrow\frac{\sqrt{2e}}{\norm V},\label{eq:angular}\\
 \frac{\sum_{k=K}^{n-1}\norm{d_k}}{\norm{z_n-z_K}}
 &\longrightarrow\sqrt{1+\frac{2e}{p^2\norm V^2}}.\label{eq:travel}
\end{align}
For $e=0$, \eqref{eq:angular} asserts $o(\log n)$, not finite length.
\end{theorem}

This is a criterion involving a finite-dimensional convex optimization profile, not a restatement of convergence of velocity directions. Neither boundedness nor convergence of the excess is asserted universally. The spectral results below establish both zero and positive excess in explicit classes.

\begin{theorem}[Regular leading spectrum]\label{thm:regular}
Let $\K$ be a finite product of real symmetric or complex Hermitian PSD cones, with arbitrary affine coupling between the blocks. If the common normalized leading state is blockwise positive definite, then $p=1/2$ and every start satisfies
\[
 z_k=V\sqrt{k}+o(\sqrt{k}),\qquad \sqrt{k}\,d_k\to V/2.
\]
There is a positive-power remainder in the velocity limit. Consequently the normalized states have finite total spherical length. Repeated nonzero eigenvalues are allowed.
\end{theorem}

The singular case is genuinely different. For $F_{ij}=E_{ij}+E_{ji}$, use indices $0,\ldots,6$ and define
\begin{align}\label{eq:family}
 a&=-\tfrac12(F_{01}+F_{34}),\\[-2mm]
 L_\gamma&=\operatorname{span}\{E_{00}+E_{33},\ E_{11}-F_{02}/2,\ E_{44}-F_{35}/2,
 \ E_{66}+\gamma(F_{02}-F_{35}),\ F_{06},\ E_{22}-E_{55}\}.\notag
\end{align}
Put
\begin{equation}\label{eq:constants}
 \begin{gathered}
 b_0=1/\sqrt3,\quad N=9/16,\quad c_0=27\sqrt3/256,
 \quad c_A=(4c_0)^{1/4},\\
 V=c_A\diag(1,0,b_0,1,0,b_0,0),\\
 F_\gamma=\frac{9(2-\sqrt3)}{64\gamma(22\gamma+1)},\qquad
 R_\gamma=\frac{N\sqrt{F_\gamma}}{c_A^3}.
 \end{gathered}
\end{equation}

\begin{theorem}[A rational family of PSD resonances]\label{thm:resonance}
For every rational $\gamma\in[3/2,5/2]$, the affine space $U_\gamma=a+L_\gamma$ is weakly infeasible for $\Ssym_+^7$ and has minimum-trace escape exponent $\alpha=1/3$. The rational matrices $Z_0(\gamma,M)$ in Section~\ref{sec:entry}, for every integer $M\geq10^{50000000000}$, generate plain DRS orbits with
\begin{align}
 Z_k&=Vk^{1/4}+Z_{\rm const}+O(k^{-1/4}),\\
 \Clust\{k^{3/4}d_k\}
 &=V/4+R_\gamma\{\cos\theta F_{06}+\sin\theta F_{26}:\theta\in\R\},\label{eq:circle}
\end{align}
where $Z_{\rm const}=-(F_{01}+F_{34})/2+(\sqrt3/4)(F_{12}+F_{45})$. Moreover
\begin{align}
 k^{3/4}\norm{d_k}_F&\longrightarrow
 \frac{\norm V_F}{4}\sqrt{1+\frac{2-\sqrt3}{\gamma(22\gamma+1)}},\label{eq:resnorm}\\
 \frac{\sum_{k=K}^{n-1}\norm{Z_{k+1}/\norm{Z_{k+1}}_F-Z_k/\norm{Z_k}_F}_F}{\log n}
 &\longrightarrow\frac23\sqrt{F_\gamma}>0.\label{eq:resangular}
\end{align}
The same conclusions hold on relatively open sets of starts in the invariant block space on $(0,1,2,6)$ and $(3,4,5)$. Block embedding gives an obstruction in every matrix order at least seven.
\end{theorem}

\paragraph{Relation to earlier work.}
The discounted-curve argument is related to the definable nonexpansive theory of Bolte--Gaubert--Vigeral \cite{BGV}, whose Theorem~3 treats the linear normalization $T^kz/k$. Liu--Ryu--Yin \cite{LRY} use divergent states and unscaled displacement limits to identify pathological conic programs. Neither statement identifies a nonzero sublinear first velocity. Do--Jentzen--Kuckuck \cite{DJK}, Theorem~1.1 in the version revised 9 November 2026, establishes finite spherical length for continuous definable gradient trajectories. Our positive and negative results concern a discrete reflected splitting; its center field need not be a convex gradient. These distinctions are theorem-level scope comparisons, not a completed priority certification.

\section{Common states, complementary shadows and the static profile}\label{sec:state}
\subsection{Discounted paths and common first terms}
For a semialgebraic nonexpansive map $T:E\to E$ with no fixed point and zero infimal displacement, fix an anchor $a$ and solve
\begin{equation}\label{eq:discount}
 z(t)=\frac{a+tTz(t)}{1+t},\qquad Tz(t)-z(t)=\frac{z(t)-a}{t}.
\end{equation}
The right-hand contraction makes $z(t)$ unique. Its graph is semialgebraic. A bounded subsequence as $t\to\infty$ would give a fixed point, so its norm diverges. Nonexpansiveness gives $\norm{z(t)-a}\leq t\norm{Ta-a}$. Comparison with an arbitrary $u$ gives
\[
 \norm{z(t)-a}\leq2\norm{u-a}+t\norm{Tu-u}.
\]
Consequently $\norm{z(t)-a}=o(t)$ when the infimal displacement is zero. One-variable Puiseux asymptotics give
\begin{equation}\label{eq:puiseux}
 z(t)-a=ct^p+O(t^{p-\eta}),\qquad
 z'(t)=pct^{p-1}+O(t^{p-1-\eta}),
\end{equation}
for $c\ne0$, rational $0<p<1$, and some $\eta>0$, reduced if necessary.

Take $t_k=t_0+k/p$ and $b_k=z(t_k)$. Equations \eqref{eq:discount}--\eqref{eq:puiseux} imply
\[
 b_{k+1}-Tb_k=O(k^{p-1-\eta})+O(k^{p-2}).
\]
Nonexpansiveness and telescoping yield $T^kz_0=b_k+o(k^p)$, and indeed a positive-power state remainder after decreasing $\eta$. Thus $V=p^{-p}c$ in \eqref{eq:state}. Two actual orbits stay a bounded distance apart. Discounted paths for different anchors do also, so $V,p$ are independent of the initial point and anchor.

For \eqref{eq:drs}, a fixed point would produce a point in $U\cap\K$. Affine points approaching the cone have displacements approaching zero. Thus the preceding argument applies under \eqref{eq:weak}. If the operator minimum-radius profile is $\Theta(\eps^{-\alpha})$, comparison with a radius-minimizing approximate fixed point and with $z(t)$ gives $p=\alpha/(1+\alpha)$. For affine-first DRS, with anchor in $U$, that operator profile is exactly
\[
 \min\{\norm{u-a}:u\in U,\ \dist(u,\K)\leq\eps\}.
\]
One inequality uses $Tu=P_\K u$ for $u\in U$; the other projects an arbitrary approximate fixed point onto $U$. For PSD constraints, a fixed right inverse of the equality map and equivalence of trace and Frobenius norm on PSD matrices identify this exponent with the minimum-trace exponent.

\subsection{The balance identities}
Moreau decomposition gives the exact actual-shadow identities
\begin{equation}\label{eq:complement}
 v_k=a+x_{k+1}+\dd y_k,\qquad
 s_k=y_{k+1}+\dd x_k,\qquad
 v_k\in\K,\ s_k\in\K^*,\ \ip{v_k}{s_k}=0.
\end{equation}
On the central path, abbreviate $x=Pz$, $y=Qz-a$ and $E(t)=\norm x^2+\norm y^2$. Its shadows are
\[
 v(t)=a+(1+1/t)x+y/t,\qquad s(t)=x/t+(1+1/t)y.
\]
Complementarity gives
\begin{equation}\label{eq:balancepath}
 \ip ay=-E(t)/t,\qquad
 \frac{t(t+1)}2E'(t)=E(t)+t\norm{y(t)}^2.
\end{equation}
For the second identity, the function $h\mapsto\ip{v(t+h)}{s(t)}$ is nonnegative with a minimum at zero, so $\ip{v'(t)}{s(t)}=0$. Substitution and differentiation of the first identity give the formula. Only the smooth central path is differentiated.

Comparing the leading Puiseux coefficients in \eqref{eq:balancepath} yields
\begin{equation}\label{eq:balance}
 H\in\K,\quad G\in\K^*,\quad \ip HG=0,\quad \ip aG=0,
 \qquad \norm G^2=p\norm V^2,\quad \norm H^2=(1-p)\norm V^2.
\end{equation}
Both $H$ and $G$ are nonzero.

\subsection{An exact optimizer and profile coefficient}
The minimum in \eqref{eq:profile} exists. A minimizing sequence has bounded $P$ component and, on a bounded residual sublevel, bounded $Q$ component. Compactness applies. Its value is positive for finite $r$, otherwise feasibility would follow, and tends to zero by \eqref{eq:weak}.

The central shadow $v(t)$ uniquely minimizes over $\K$ the strictly convex coercive function
\begin{equation}\label{eq:centralopt}
 F_t(v)=\frac{\norm{Pv}^2}{2(t+1)^2}+\frac{\norm{Qv-a}^2}{2}.
\end{equation}
Its gradient at $v(t)$ is $s(t)/(t+1)\in\K^*$ and complementary to $v(t)$, which is precisely the optimality inequality. Hence
\[
 D(\norm{Pv(t)})=\norm{Qv(t)-a}.
\]
The radius is eventually strictly increasing by its Puiseux expansion and covers a tail. Eliminating $t$ using \eqref{eq:puiseux} gives
\begin{equation}\label{eq:Dconstant}
 D(r)\sim p\norm G\norm H^{(1-p)/p}r^{-(1-p)/p},
 \qquad D(r_k)^2\sim p^2\norm G^2k^{2p-2}.
\end{equation}
This is an exact leading coefficient, not merely a power comparison.

\section{Finite-window energy and the universal defect law}
\subsection{Exact two-shadow identities}
Put $q_k=\norm{d_k}^2$, $w_k=\ip{y_{k+1}}{\dd y_k}$ and
\[
 \phi_k=\ip{v_{k+1}}{s_k}\geq0,\qquad
 \psi_k=\ip{v_k}{s_{k+1}}\geq0.
\]
Firm nonexpansiveness and \eqref{eq:complement} imply
\begin{align}
 q_k-q_{k+1}&\geq\norm{d_{k+1}-d_k}^2,\label{eq:qmono}\\
 \norm{\dd x_{k+1}}^2
 &=w_k-w_{k+1}+\ip{\dd y_k}{\dd y_{k+1}}-\psi_k,\label{eq:twoshadow}\\
 \phi_k+\psi_k&=\tfrac12\bigl(q_k-q_{k+1}-\norm{d_{k+1}-d_k}^2\bigr).\label{eq:cross}
\end{align}
For \eqref{eq:cross}, expand $\ip{v_{k+1}-v_k}{s_{k+1}-s_k}$ in two ways; complementarity makes it $-(\phi_k+\psi_k)$. In particular $\sum_{j=n}^{m-1}\psi_j\leq q_n/2$. Polarizing the normal inner product in \eqref{eq:twoshadow}, summing and shifting one endpoint gives
\begin{equation}\label{eq:exactenergy}
 \sum_{j=n}^{m-1}q_j=2\sum_{j=n}^{m-1}\delta_j^2+w_n-w_m+R_{n,m},
 \qquad |R_{n,m}|\leq4q_n.
\end{equation}
All these are finite identities and inequalities, also when $p\leq1/2$; no divergent energy tail is used.

\subsection{Natural scale and radial rigidity}
If $X_n(\lambda)$ is bounded for one fixed ratio, covering finitely many such windows bounds $\sum\delta_j^2=O(n^{2p-1})$ on every fixed-ratio interval. Choose indices $i\in[n/4,n/3]$ and $l\in[2n,3n]$ for which $\delta_i,\delta_l=O(n^{p-1})$, by averaging. Since $y_j=O(j^p)$, their $w$ values are $O(n^{2p-1})$. Summing \eqref{eq:twoshadow} between those endpoints and discarding $-\psi$ bounds the tangential squared energy by $O(n^{2p-1})$. The normal energy has the same bound. Monotonicity of $q_k$ then yields $q_n=O(n^{2p-2})$. The converse follows immediately from \eqref{eq:Dconstant} and $e_k\leq q_k$.

For any nonexpansive orbit with \eqref{eq:state} and $d_k=O(k^{p-1})$, monotonicity of $I-T$ gives
\[
 \ip{d_k-d_j}{z_k-z_j}\leq0.
\]
For $j/k$ in a compact interval separated from one, substitute the state expansion and the step bound. Average over $j\in[(1+h)k,(1+2h)k]$, and separately over $[(1-2h)k,(1-h)k]$. The sums of $d_j$ telescope. Sending $k\to\infty$ and then $h\downarrow0$ squeezes $k^{1-p}\ip{d_k}{V}$ to $p\norm V^2$. This proves \eqref{eq:radial} without differentiating the state remainder.

Cross-pair $v_k$ with $s(t)$ and $v(t)$ with $s_k$ at $t=k/p$. Their pairings are nonnegative. Substitution of the central and actual states, \eqref{eq:balancepath}, and the natural step bound gives
\[
 \liminf k^{1-p}\ip{\dd y_k}{G}\geq p\norm G^2,\qquad
 \liminf k^{1-p}\ip{\dd x_k}{H}\geq p\norm H^2.
\]
For the second cross-pairing, use additionally the exact actual identity
\[\ip{a}{y_{k+1}}=-\ip{x_{k+1}}{\dd x_k}-\ip{y_{k+1}}{\dd y_k}.\]
It follows from complementarity; its leading coefficient follows from the total radial limit. The sum of the two lower bounds is that same radial limit, so both are limits and equalities. Therefore
\begin{equation}\label{eq:wlimit}
 w_k=p\norm G^2 k^{2p-1}+o(k^{2p-1}).
\end{equation}

\subsection{Proof of Theorem~\ref{thm:excess}}
Insert \eqref{eq:Dconstant}, \eqref{eq:wlimit} and $q_n=O(n^{2p-2})$ into \eqref{eq:exactenergy}. The coefficient of the baseline energy becomes
\[
 2p^2\norm G^2-p(2p-1)\norm G^2=p\norm G^2=p^2\norm V^2.
\]
At $p=1/2$, the leading $w$ endpoints cancel, while $I_{n,m}\to\log\lambda$; the same formula holds. This proves \eqref{eq:energy}. Expanding the squared distance and using \eqref{eq:radial} proves \eqref{eq:transverse}.

Suppose $X_n(\lambda)\to0$ for one fixed $\lambda$. The radial limit implies the uniform tail lower bound
\[
 q_j\geq[A_0-o(1)]j^{2p-2}.
\]
If $q_n\geq(A_0+\eps)n^{2p-2}$ along a subsequence, monotonicity of $q$ preserves a fixed positive excess over the baseline on a sufficiently short proportional interval immediately before $n$. It lies inside a window of the prescribed ratio ending at $n$; the uniform lower bound elsewhere prevents cancellation. This contradicts \eqref{eq:energy}. Hence $k^{2-2p}q_k\to A_0$, and \eqref{eq:radial} identifies the vector as $pV$. The converse is immediate. Pointwise saturation implies $e_k=o(k^{2p-2})$, hence the window condition; the vector law implies saturation by applying $Q$ and using \eqref{eq:Dconstant}.

For nonzero $e$, the energy formula for every ratio and monotone-density secants yield \eqref{eq:normlimit}. Conversely its Riemann sums and the energy formula give $X_n(\lambda)\to e$ for every ratio. Radial normalization has the finite-difference expansion
\[
 k(\nu_{k+1}-\nu_k)=
 \frac{(I-\nu_k\nu_k^*)k^{1-p}d_k}{\norm{z_k}/k^p}+O(k^{-1}).
\]
The natural-scale bound, radial rigidity and state limit give
\[
 k^2\norm{\nu_{k+1}-\nu_k}^2
 =\frac{\norm{k^{1-p}d_k-pV}^2}{\norm V^2}+o(1).
\]
Under \eqref{eq:normlimit} the right side tends to $2e/\norm V^2$. Summation proves \eqref{eq:angular}. Summing $\norm{d_k}\sim\sqrt{A_0+2e}\,k^{p-1}$ proves \eqref{eq:travel}. This completes the theorem.

\paragraph{Residual descriptions.}
For any equality representation $U=\{z:\mathcal A z=b\}$, including redundant equations,
\[
 \mathcal A v_k-b=\mathcal A\dd y_k,\qquad
 \norm{(\mathcal A\mathcal A^*)^{\dagger/2}(\mathcal A v_k-b)}=\delta_k.
\]
The identity follows from $\mathcal A^*(\mathcal A\mathcal A^*)^\dagger\mathcal A=Q$. Under zero excess the raw vector limit is $p\mathcal A G\ne0$. Nonvanishing follows from $G\in L^\perp\setminus\{0\}$. Ill-conditioning can change its magnitude, but redundancy cannot make it vanish.

\section{Why a convex splitting has no leading kinetic defect}
The following comparison result records the convex mechanism separately. It is not a change of algorithm in Theorem~\ref{thm:resonance}.

\begin{theorem}[Convex forward--backward contrast]\label{thm:fb}
Let $F=g+h$ be proper closed convex and semialgebraic, with finite unattained infimum $F_*$. Let $g$ be finite convex with globally $L$-Lipschitz gradient and $h$ proper closed convex. Neither summand separately needs to be definable. For fixed $0<\eta<2/L$ (any $\eta>0$ if $L=0$), iterate
\[
 x_{k+1}=\operatorname{prox}_{\eta h}(x_k-\eta\nabla g(x_k)),\qquad\tau_k=\eta k.
\]
There are a nonzero $V$ and rational $0<p<1/2$, determined by $F$ and the metric and not by its splitting, such that every start satisfies
\begin{align*}
 x_k&=V\tau_k^p+o(\tau_k^p),&
 \tau_k^{1-p}\frac{x_{k+1}-x_k}{\eta}&\longrightarrow pV,\\
 F(x_k)-F_*&\sim\frac{p^2\norm V^2}{1-2p}\tau_k^{2p-1},&
 \frac{F(x_k)-F(x_{k+1})}{\eta}&\sim p^2\norm V^2\tau_k^{2p-2}.
\end{align*}
The normalized states have finite total spherical length. The same state and velocity coefficients apply to any admissible convex splitting at equal physical time.
\end{theorem}

\begin{proof}
Let $y(s)$ uniquely minimize $F(y)+\norm{y-a}^2/(2s)$ and set $u(s)=(a-y(s))/s\in\partial F(y(s))$. Comparing to arbitrary approximate minimizers gives $F(y(s))\to F_*$ and $\norm{y(s)-a}^2/s\to0$. Nonattainment forces divergence. Puiseux asymptotics give
\[
 y(s)-a=cs^p+o(s^p),\qquad y'(s)=pcs^{p-1}+o(s^{p-1}),\qquad 0<p<1/2.
\]
The subgradient inequalities along this smooth path give
\[
 \frac{d}{ds}F(y(s))=\ip{u(s)}{y'(s)}
 =-p\norm c^2s^{2p-2}+o(s^{2p-2}).
\]
Integration to infinity determines the gap coefficient. Put $V=p^{-p}c$ and $B=p^2\norm V^2/(1-2p)$.

The total objective, rather than its chosen splitting, supplies this central path. To track it by the possibly nontame splitting, write $b(s)=\nabla g(y(s))$. Cocoercivity of $\nabla g$ and monotonicity of $\partial h$ imply
\[
 \norm{b(s+v)-b(s)}^2
 \leq L\ip{u(s+v)-u(s)}{y(s+v)-y(s)}
 =O(v^2s^{2p-3})
\]
for bounded positive $v$. Thus the gradient-splitting correction is of smaller order than $s^{p-1}$. Comparing the proximal input for $y(s_{k+1})$ with the forward--backward input for $y(s_k)$, at $s_k=s_0+\eta k/p$, gives local error $o(\eta s_k^{p-1})$. Nonexpansiveness of the update and summation prove the common state term.

Here are the energy steps that actually determine the velocity. Subtract $F_*$ and put
\[
 F_k=F(x_k)-F_*,\quad G_k=(x_k-x_{k+1})/\eta,\quad
 q_k=\norm{G_k}^2,\quad e_k=\nabla g(x_{k+1})-\nabla g(x_k).
\]
Proximal optimality gives $G_k+e_k\in\partial F(x_{k+1})$. With $d_\eta=1-\eta L/2>0$, the descent and comparator inequalities are
\begin{align*}
 F_k-F_{k+1}&\geq\eta d_\eta q_k,\\
 2\eta[F_{k+1}-(F(w)-F_*)]
 &\leq\norm{x_k-w}^2-\norm{x_{k+1}-w}^2
 -(1-\eta L)\norm{x_{k+1}-x_k}^2.
\end{align*}
When $\eta L>1$, the positive remainder in the second inequality is controlled by summing the first. Therefore, after the first step, a fixed finite $C$ gives
\[
 F_k\leq F(w)-F_*+\frac{\norm{x_1-w}^2+C}{2\eta(k-1)}.
\]
Choose $w=y(\tau_k)$. This gives $F_k=O(\tau_k^{2p-1})$. Nonexpansiveness makes $q_k$ nonincreasing; summing descent from $k/2$ to $k$ yields $q_k=O(\tau_k^{2p-2})$. The actual subgradient $G_{k-1}+e_{k-1}$ has the same $O(\tau_k^{p-1})$ bound. Comparing it and $u(s_k)$ against the already proved $o(\tau_k^p)$ state difference squeezes
\[
 F_k-F(y(s_k))+F_*=o(\tau_k^{2p-1}),\qquad F_k\sim B\tau_k^{2p-1}.
\]

The forward contraction, followed by nonexpansiveness of the proximal map, gives for $L>0$
\[
 \sum_{j\geq k}\norm{e_j}^2\leq C_\eta q_k,
 \qquad C_\eta=\frac{\eta}{2/L-\eta}.
\]
For $L=0$ these errors vanish. The consecutive subgradient inequalities imply
\[
 \eta(q_j+\ip{e_j}{G_j})\leq F_j-F_{j+1}
 \leq\eta(\ip{G_{j-1}}{G_j}+\ip{e_{j-1}}{G_j}).
\]
Let $I_k=\eta\sum_{j\geq k}q_j<\infty$. Summation, polarization, monotonicity of $q$ and Cauchy--Schwarz yield
\[
 |F_k-I_k|\leq\eta q_{k-1}+2\sqrt{\eta C_\eta q_{k-1}I_k}.
\]
Since $I_k\leq F_k/d_\eta$, this error is $o(F_k)$. Consequently
\[
 \eta\sum_{j\geq k}q_j\sim B(\eta k)^{2p-1}.
\]
Forward and backward fixed-ratio secants of this tail, followed by a ratio tending to one, give
\[
 q_k\sim (1-2p)B\tau_k^{2p-2}=p^2\norm V^2\tau_k^{2p-2}.
\]
The nonexpansive radial-rigidity argument in Section 3 identifies the velocity vector. This is the missing step that would be invalid if replaced by differentiation of $x_k=V\tau_k^p+o(\tau_k^p)$. Moreover $q_{k+1}/q_k\to1$ and the one-step contraction give $e_k=o(\tau_k^{p-1})$. The displayed subgradient bounds now squeeze the individual objective decrement to the stated coefficient.

For completeness, finite angular length requires a rate, not mere direction convergence. Choose $0<\kappa<p<1/2$ so that the Puiseux path and its derivative have relative $O(s^{-\kappa})$ remainders. Tracking gives $x_k=V\tau_k^p+O(\tau_k^{p-\kappa})$ and the subgradient comparison gives $F_k=B\tau_k^{2p-1}[1+O(\tau_k^{-\kappa})]$. The energy-tail error above has smaller relative order $O(\tau_k^{-1/2})$. Monotone-density secants of relative width $\tau_k^{-\kappa/2}$ therefore give a relative $O(\tau_k^{-\kappa/2})$ remainder for $q_k$. The one-step contraction bounds $e_k$ by $O(\tau_k^{p-1-\kappa/4})$. Convex secants between $x_k$ and $x_{\lfloor(1\pm\tau_k^{-\kappa/2})k\rfloor}$, using the state and gap remainders, determine the radial pairing to error $O(\tau_k^{-\kappa/2})$. Combining norms and pairings gives a velocity remainder $O(\tau_k^{-\kappa/8})$ after rescaling. Radial normalization then has increments $O(k^{-1-\kappa/8})+O(k^{-2})$, which are summable.
\end{proof}

For alternating projections take $g(x)=\dist(x,C)^2/2$ and $h=\iota_U$. Theorem~\ref{thm:fb} then covers semialgebraic zero-distance disjoint convex pairs. This is a consequence for a different algorithm, not a proof of the DRS law. The decisive convex ingredient is the asymptotically exact objective/energy tail. In conic DRS the finite-window identity instead contains the nonnegative static excess, which the rational resonance makes strictly positive.

\section{Exact matrix transport and regular leading spectra}
\subsection{A cut identity}
For a PSD cone, write $u=P_UZ$, $v=(2u-Z)_+$, $s=(Z-2u)_+$ and $d=v-u$. Then $v,s\succeq0$, $vs=sv=0$, $v=u+d$ and $s=Z-u+d$.

Suppose $U$ is invariant under congruence by $J_c=\diag(I,-I)$. The affine offset is block diagonal and the affine off-diagonal projection has the form $u_{12}=\Pi W$, where $W=Z_{12}$ and $\Pi$ is an orthogonal projector. Thus
\[
 v_{12}=\Pi W+\dd W,\qquad s_{12}=(I-\Pi)W+\dd W.
\]
Taking the trace of the $(2,2)$ block of $vs=0$ gives
\begin{equation}\label{eq:cut}
 \norm{W+\dd W}_F^2+\norm{\dd W}_F^2
 +2\tr(v_{22}s_{22})=\norm W_F^2.
\end{equation}
The overlap is nonnegative and equals $\tr(v_{11}s_{11})$. Consequently both squared transverse increments and overlap are summable, with their exact tail sum equal to the lost cut energy.

For a rank-one cut $R=\left[\begin{smallmatrix}A&b\\b^*&c\end{smallmatrix}\right]$, there is also an exact transport representation. In an eigenbasis of $A$, let $\lambda_i$ be its eigenvalues and let $(\mu_j,q_j)$ be spectral data of $R$. Define
\[
 C=g(A),\qquad g(\lambda)=\sum_j |(q_j)_{\rm last}|^2
 \frac{|\mu_j|-|\lambda|}{\mu_j-\lambda},
\]
with the continuous divided difference at equal nonzero arguments and any value in $[-1,1]$ at $(0,0)$. Then
\begin{equation}\label{eq:transport}
 (|R|)_{12}=Cb,\quad -I\preceq C\preceq I,\quad
 \ip b{(I-C^2)b}=4v_{\rm last,last}s_{\rm last,last}.
\end{equation}
To verify the first formula, use $(\mu_j-\lambda_i)(q_j)_i=b_i(q_j)_{\rm last}$ and subtract $|\lambda_i|$ times the spectral orthogonality identity. Repeated eigenspaces cause no ambiguity. The last formula is the bottom diagonal of $|R|^2=R^2$. In DRS the reflected coupling is $(2\Pi-I)W$, so
\[
 W_{\rm next}=\tfrac12[I+C(2\Pi-I)]W.
\]
This is the current nonlinear operator, not its moving linearization.

\subsection{Proof of the regular-spectrum theorem}
We give the argument blockwise, allowing the affine projection to couple blocks. Write $J=2P-I$, so $TZ=(Z+|JZ+2a|)/2$. Let the unit cosmic limit $N$ be blockwise positive definite. Its complementary parts $H=PN$, $G=QN$ have orthogonal support projections $E,F$, with $E+F=I$ in each block. Define
\[
 D_H=\{X\in L:X=EXE\},\quad
 D_G=\{Y\in L^\perp:Y=FYF\},\quad D=D_H\oplus D_G.
\]
Every $Y\in L^\perp\cap\K$ has $\ip aY\geq0$ by weak feasibility. Complementarity and the radial identity below imply $\ip aG=0$; equivalently this follows from \eqref{eq:balance}. Applying the inequality to $G\pm\eps Y$ for $Y\in D_G$ gives $a\perp D$.

For $n\in D$ near $N$, the matrix $Jn$ is nonsingular and the absolute-value map is analytic there. Let $C_n=D|\cdot|(Jn)$ and $M_n=(I+C_nJ)/2$. In eigenbases of $n=H_n+G_n$, the derivative multiplies entries by $1$ on $E$--$E$, by $-1$ on $F$--$F$, and by
\[
 \frac{h_i-g_j}{h_i+g_j}\in(-1,1)
\]
on the cross blocks. Thus $M_n$ is firmly nonexpansive, fixes precisely $D$, preserves $D^\perp$, and satisfies $\norm{M_n|_{D^\perp}}\leq q<1$ uniformly near $N$. Indeed equality in its norm contraction would, by firmness, give a fixed vector; equality in the cross coefficient bound forces support-preserving components. This identifies the fixed space exactly.

Decompose $Z=s+w$ with $s\in D$, $w\perp D$, $r=\norm s$, $n=s/r$. Cosmic convergence gives $w=o(r)$. Uniform Taylor expansion yields
\begin{align}\label{eq:regularTaylor}
 T(s+w)&=s+M_nw+C_na+r^{-1}F_1(n,w)
 +O((\norm w+1)^3/r^2),\\
 F_1(n,w)&=\tfrac14D^2|\cdot|(Jn)[Jw+2a,Jw+2a].\notag
\end{align}
The first-order remainder and strict contraction first show $w=O(1)$. Projection on $D$ gives $\dd s=O(r^{-1})$, $r_{\rm next}/r=1+O(r^{-2})$, and $\dd n=O(r^{-2})$. Define
\[
 w_0(n)=(I-M_n|_{D^\perp})^{-1}C_na,
 \quad w_1(n)=(I-M_n|_{D^\perp})^{-1}\Pi_{D^\perp}F_1(n,w_0(n)).
\]
Subtract these targets successively in \eqref{eq:regularTaylor}. Rescaling the contractive recurrences first by $r$ and then by $r^2$ proves
\[
 w=w_0(n)+r^{-1}w_1(n)+O(r^{-2}),\qquad
 d=r^{-1}b(n)+O(r^{-2}),\quad b(n)=\Pi_DF_1(n,w_0(n)).
\]
Differences of two $O(r^{-2})$ terms remain $O(r^{-2})$; no differentiation of a discrete remainder is needed.

Scalar complementarity gives the exact radial relation
\begin{equation}\label{eq:radialexact}
 \norm{TZ}^2-\norm Z^2=2h(Z)-\norm d^2,
 \qquad h(Z)=\norm a^2-\ip aZ.
\end{equation}
At the stationary normal point set $B=Jw_0+2a$. The equation $w_0=C_nB$ yields
\[
 4h_0(n):=4(\norm a^2-\ip a{w_0})=\norm B^2-\norm{C_nB}^2\geq0.
\]
Equality would force support-block diagonal $B,w_0$. Then $a+Pw_0$ is supported on $E$, and adding a sufficiently large multiple of the positive recession $H_n$ makes it PSD, contradicting infeasibility. Hence $h_0(N)>0$.

Equation \eqref{eq:radialexact} now gives $r_k/\sqrt{k}\to\sqrt{2h_0(N)}$ and $\sqrt{k}d_k\to b(N)/\sqrt{2h_0(N)}$. Summing this actual step limit and comparing with the cosmic direction yields $b(N)=h_0(N)N$. Thus $V=\sqrt{2h_0(N)}N$ and the limit is $V/2$. Every orbit has the same cosmic limit by nonexpansiveness, so the result holds for all starts.

The positive-power state remainder from Section~\ref{sec:state}, smoothness of $b$ and the displayed tracking estimates give $\sqrt{k}d_k=V/2+O(k^{-\zeta})$ for some $\zeta>0$. Radial normalization cancels the leading radial velocity; its increments are $O(k^{-1-\zeta})+O(k^{-2})$, proving finite spherical length.

\section{The rational resonance: geometry and exact local equations}
\subsection{Weak infeasibility and the rational starting family}\label{sec:entry}
An equality description of \eqref{eq:family} is
\[
 X_{00}=X_{33},\quad X_{22}+X_{55}=0,\quad
 X_{11}+2X_{02}-2\gamma X_{66}=0,\quad
 X_{44}+2X_{35}+2\gamma X_{66}=0,
\]
with $X_{01}=X_{34}=-1/2$ and all other off-diagonal pairs zero except $(0,2),(0,6),(3,5)$. In a PSD solution, $X_{22}=X_{55}=0$ kills their rows. The two displayed equations then force $X_{11}=2\gamma X_{66}$ and $X_{44}=-2\gamma X_{66}$; positivity kills all three diagonals and contradicts $X_{01}=-1/2$.

For $u>0$, two identical blocks
\[
 \begin{pmatrix}(2u)^{-1}&-1/2&-u/2\\-1/2&u&0\\-u/2&0&u^3\end{pmatrix}
\]
and a zero seventh row are PSD by the leaf Schur complement. Their only residual is $X_{22}+X_{55}=2u^3$, and their trace is of order $u^{-1}$. Conversely, a PSD point of equality residual at most $\delta$ has $X_{55}\leq\delta$, $|X_{34}|\geq1/4$ for small $\delta$, and, writing $B=X_{33}$, $d=X_{44}$, $r=X_{66}$,
\[
 d+2\gamma r\leq2\sqrt{B\delta}+\delta,
 \qquad 1/16\leq Bd\leq2B^{3/2}\sqrt\delta+B\delta.
\]
Thus $\tr X\geq c\delta^{-1/3}$, proving the exact exponent.

For an integer $M$ set $t=M^{-1}$ and $s=\operatorname{isqrt}(3M^8)/M^4$, where $\operatorname{isqrt}$ is the floor of the nonnegative square root. Put
\[
 J_0=-s(14\gamma+1)/30000,\quad R_0=3/2500,\quad
 C_f=2sJ_0+2\gamma R_0,
\]
\[
 C_a=C_f+3J_0/2,\quad G_a=\gamma R_0-C_f/2+3J_0/2,\quad
 Q_0=R_0-6\gamma J_0.
\]
Define the symmetric $Z_0(\gamma,M)$ by
\begin{align*}
 Z_{00}=Z_{33}&=M,& Z_{22}=Z_{55}&=sM/3,\\
 Z_{11},Z_{44}&=t(3/8+3s/16)\ \pm t^3C_a,&
 Z_{02},Z_{35}&=3t(s-1)/16\ \pm t^3G_a,\\
 Z_{01}=Z_{34}&=-1/2+9st^2/64,&
 Z_{12},Z_{45}&=s/4\ \pm t^2J_0,\\
 Z_{06}&=t/50,& Z_{16}=Z_{26}&=0,\qquad Z_{66}=t^3Q_0,
\end{align*}
and all other entries zero. Every entry is rational when $\gamma$ is rational. The two blocks $(0,1,2,6)$ and $(3,4,5)$ are invariant under both projections.

\subsection{Weighted variables and the exact affine projection}
For a state in this block space use
\begin{align*}
 A&=(Z_{00}+Z_{33})/2,& B&=(Z_{22}+Z_{55})/2,& t&=A^{-1},\quad\tau=t^2,\quad b=B/A,\\
 H&=(Z_{12}+Z_{45})/2,& j&=(Z_{12}-Z_{45})/(2t^2),\\
 \xi&=Z_{06}/t,&\eta&=Z_{26}/t,&\chi&=Z_{16}/t^2,\\
 C&=(Z_{11}+Z_{44})/(2t),&G&=(Z_{02}+Z_{35})/(2t),\\
 C_a&=(Z_{11}-Z_{44})/(2t^3),&G_a&=(Z_{02}-Z_{35})/(2t^3),&Q_0&=Z_{66}/t^3.
\end{align*}
Also set $A_d=(Z_{00}-Z_{33})/2=t^5a_d$, $B_d=(Z_{22}-Z_{55})/2=t^5b_d$, $(Z_{01}-Z_{34})/2=t^4f_a$, and $(Z_{01}+Z_{34})/2=-1/2+t^2f$. The exact Gram equations for the affine projection are
\begin{equation}\label{eq:affinegamma}
 c=\frac{2(C-G)}3,\quad r=\frac{3Q_0+4\gamma C_a+8\gamma G_a}{3+8\gamma^2},\quad
 c_a=\frac{2(C_a-G_a+\gamma r)}3,\quad g_a=-c_a/2+\gamma r.
\end{equation}
In physical coordinates $u_{11},u_{44}=tc\pm t^3c_a$, $u_{02},u_{35}=-tc/2\pm t^3g_a$, $u_{66}=t^3r$, $u_{00}=u_{33}=A$, $u_{22}=-u_{55}=B_d$, $u_{06}=Z_{06}$ and $u_{01}=u_{34}=-1/2$.

For $D_0=\diag(1,-1,-3,1,-1,-3,-1)$, the anticommutator $X\mapsto D_0X+XD_0$ annihilates $a$ and preserves $L_\gamma$. Matrix complementarity therefore proves the exact charge identity
\begin{equation}\label{eq:charge}
 E=\tr(D_0Z^2),\qquad E_{\rm next}-E=-\tr(D_0d^2).
\end{equation}
Its coordinate expression determines $b$ analytically near $b_0$:
\begin{align}\label{eq:bcharge}
 b^2={}&\tfrac13-\tfrac\tau6(E+8H^2)
 -\tfrac{\tau^2}6(2C^2+4G^2+4\eta^2)
 -\tfrac{\tau^3}6(8j^2+2\chi^2)\\
 &-\tfrac{\tau^4}6(2C_a^2+4G_a^2+Q_0^2)
 +\tfrac{\tau^6}6(2a_d^2-6b_d^2).\notag
\end{align}

\subsection{An analytic chart for the actual PSD projection}\label{sec:chart}
In coordinates $(0,1;2,6)$ the reflected block $R=2u-Z$ is
\begin{equation}\label{eq:Rblocks}
 R_{PP}=\begin{pmatrix}a/t&-\beta\\-\beta&td\end{pmatrix},\quad
 R_{QQ}=\begin{pmatrix}-B_r/t&-t\eta\\-t\eta&t^3q\end{pmatrix},\quad
 R_{QP}=\begin{pmatrix}-tK&-h\\t\xi&-t^2\chi\end{pmatrix},
\end{equation}
where
\begin{gather*}
 a=1-\tau^3a_d,\ B_r=b-\tau^3b_d,\ \beta=1/2+\tau f+\tau^2f_a,\\
 d=2c-C+\tau(2c_a-C_a),\quad h=H+\tau j,\quad
 K=c+G-\tau(2g_a-G_a),\quad q=2r-Q_0.
\end{gather*}
Represent the positive invariant subspace as the graph
\[
 X=\begin{pmatrix}t^2x&ty\\t^2z&t\omega\end{pmatrix}.
\]
Dividing the entries of $R_{QP}+R_{QQ}X-XR_{PP}-XR_{PQ}X=0$ by $t,1,t,t^2$, respectively, gives the exact equations
\begin{align}\label{eq:riccati}
0={}&-K-(a+B_r)x+\beta y-\tau\eta z+\tau xyh+\tau^2x^2K-\tau^2z(x\xi-y\chi),\\
0={}&-h-B_ry+\tau(-\eta\omega+\beta x-dy+y^2h)
 +\tau^2(xyK-\omega(x\xi-y\chi)),\notag\\
0={}&\xi-az+\beta\omega+\tau(-\eta x+x\omega h)
 +\tau^2(qz+xzK-z(z\xi-\omega\chi)),\notag\\
0={}&-\chi-\eta y+\beta z-d\omega+y\omega h
 +\tau(q\omega+yzK-\omega(z\xi-\omega\chi)).\notag
\end{align}
The second block has opposite asymmetric signs and no $z,\omega,\xi,\eta,\chi$ coupling. At $\tau=0$ the solutions are
\[
 y=-H/b_0,\quad x=(-c-G+y/2)/(1+b_0),\quad
 \delta=2c-C+H^2/b_0-1/4,
\]
\[
 \omega=(\xi/2-\chi+H\eta/b_0)/\delta,\qquad z=\xi+\omega/2.
\]
At the symmetric reference values
\[
 H_0=\sqrt3/4,\quad C_*=3/8+3\sqrt3/16,\quad
 G_*=3(\sqrt3-1)/16,\quad\chi=3\eta/4,
\]
one has $\delta=1/8$, $x=-N$, $y=-3/4$, $z=3\xi$, $\omega=4\xi$. The combined Jacobian of the charge and both graph systems has determinant, up to sign, $2b_0^3(1+b_0)^2\delta$. It is uniformly invertible on a sufficiently small compact neighborhood.

The sign of the remaining small Schur eigenvalue will be controlled below. Once negative, the analytic graph is the actual positive spectral subspace. Its two positive eigenvalues have a positive middle Schur coefficient $t/8+O(t^3)$. The other block has two positive eigenvalues and one negative eigenvalue. Repeated eigenvalues within a sign subspace do not affect the graph construction.

\section{Fast contraction, nonlinear averaging and all-time trapping}
\subsection{Leading recurrences and their exact coefficients}
Define $\Psi(H,b)$ by
\[
 c_f(H,b)=2H^2/b^2-H/b,\qquad
 \Psi(H,b)=(c_f+H^2/b,\ H^2/b-c_f/2).
\]
The symmetric fast matrix at the reference is
\[
 M_s=\begin{pmatrix}2/3&-2/3\\(-2+\sqrt3)/3&(1+\sqrt3)/6\end{pmatrix},
\]
strictly contractive for $C^2+2G^2$. The slow symmetric filter is
\[
 \dd H=\tau\frac{H^2(3b/2-2H)}{b^3}
 +O\bigl(\tau\norm{(C,G)-\Psi(H,b)}+\tau^2\bigr).
\]
Its derivative at $(H_0,b_0)$ is $-\kappa$, where $\kappa=9\sqrt3/8$.

For the asymmetric fast coordinates the leading update, after symmetric tracking, is
\begin{equation}\label{eq:fast}
 C_a'=c_a+3j/2,\quad
 G_a'=\frac{b_0G_a+(1-b_0)g_a+(j-\xi\eta)/2}{1+b_0},\quad
 Q_0'=Q_0+3\xi^2-r,
\end{equation}
up to $O(\tau)$. Its fixed point is
\begin{equation}\label{eq:psia}
 \begin{gathered}
 r_0=3\xi^2,\qquad c_f^a=2\sqrt3j+2\gamma r_0+\sqrt3\xi\eta,\\
 \psi_\gamma(j,\xi,\eta)=
 (c_f^a+3j/2,\ \gamma r_0-c_f^a/2+3j/2,\ r_0-6\gamma j).
 \end{gathered}
\end{equation}
Let $M_\gamma$ denote its linear part and $W=\diag(2,4,1)$. Direct multiplication gives the three leading principal minors of $W-M_\gamma^TWM_\gamma$:
\begin{equation}\label{eq:minors}
 \frac{16\gamma^2-18+16\sqrt3}{(3+8\gamma^2)^2},\quad
 \frac{(-128+96\sqrt3)\gamma^2-64+44\sqrt3}{(3+8\gamma^2)^2},\quad
 \frac{-64+44\sqrt3}{(3+8\gamma^2)^2}.
\end{equation}
They are positive for $\gamma\geq0$. On $[3/2,5/2]$ the weighted norm contracts by at most $1-10^{-4}$; one obtains this from the normalized defect determinant, which exceeds $1/2000$, and the fact that its eigenvalues are at most one.

After the fast errors are $O(\tau)$, substitution into \eqref{eq:riccati} or into $vs=0$ gives
\begin{align}\label{eq:phasefilter}
 \dd\xi&=N\tau\eta+O(\tau^2),&\dd\eta&=-N\tau\xi+O(\tau^2),\\
 \dd j&=\tau[-\kappa j-(9\gamma/2)\xi^2-(3\sqrt3/4)\xi\eta]+O(\tau^2).\notag
\end{align}
These coefficients retain the nonlinear diagonal feedback. For example, before setting $r=3\xi^2$, the antisymmetric $12$ complementarity equation gives
\[
 \frac1{b_0}\bigl[(3/4)j-(3/8)j-H_0c_a\bigr]
 =-\kappa j-(3\gamma/2)r-(3\sqrt3/4)\xi\eta.
\]

Put
\begin{equation}\label{eq:Sgamma}
 S_\gamma=\frac{\sqrt3}2\gamma(14\gamma+1)\xi^2
 +\frac{\sqrt3}2\gamma(2\gamma-1)\eta^2
 -3\gamma(2\gamma-1)\xi\eta.
\end{equation}
The periodic target $j_\gamma=-S_\gamma/(6\gamma)$ satisfies exactly
\[
 N(\eta\partial_\xi-\xi\partial_\eta)j_\gamma
 =-\kappa j_\gamma-(9\gamma/2)\xi^2-(3\sqrt3/4)\xi\eta.
\]
At these targets,
\begin{equation}\label{eq:radialamplitudecoeff}
 v_{66}=3\xi^2t^3+O(t^5),\quad s_{66}=S_\gamma t^3+O(t^5),
 \quad\dd A=c_0t^3+O(t^5),\quad t'=t-c_0t^5+O(t^7).
\end{equation}
The radial coefficients follow from the positive and negative Schur complements: $v_{22}=(27/256)t^3+O(t^5)$ and $s_{00}=b_0N^2t^3+O(t^5)=c_0t^3+O(t^5)$.

The last small Schur coefficient is $-(S_\gamma-\sqrt3\eta^2)$. The determinant of this quadratic form is
\begin{equation}\label{eq:sign}
 \det(S_\gamma-\sqrt3\eta^2)=\frac{3\gamma}{2}(8\gamma^3-16\gamma-1).
\end{equation}
Its first diagonal coefficient is positive. It is therefore positive definite exactly above the positive root $\gamma_c$ of $8\gamma^3-16\gamma-1=0$. On $[3/2,5/2]$, its determinant is at least $9/2$ and trace is below $100$, giving least eigenvalue above $9/200$. This supplies a uniform sign margin on an annulus.

\subsection{The corrected amplitude}
Write $S_\gamma=A_\gamma\xi^2+B_{2,\gamma}\eta^2+C_\gamma\xi\eta$ and put
\[
 B_\gamma=\frac{3\sqrt3}{4}\gamma(22\gamma+1),\qquad
 a_0=2c_0-N^2=\frac{27(2\sqrt3-3)}{256}>0.
\]
Define the quartic polynomial
\begin{align}\label{eq:Pgamma}
 P_\gamma={}&\frac{C_\gamma}{N}
 \left(\frac{15}{16}\xi^4-\frac98\xi^2\eta^2-\frac9{16}\eta^4\right)\\
 &+\frac{B_\gamma-6A_\gamma}{N}\xi^3\eta
 -\frac{B_\gamma}{N}\xi\eta^3.\notag
\end{align}
Exact differentiation proves
\[
 N(\eta\partial_\xi-\xi\partial_\eta)P_\gamma
 =6\xi^2S_\gamma-B_\gamma(\xi^2+\eta^2)^2.
\]
For $W_c=(Z_{06},Z_{16},Z_{26})$, let
\[
 F=A^2\norm{W_c}^2=\xi^2+\eta^2+\tau\chi^2,
 \qquad \widetilde F=F+\tau P_\gamma.
\]
The exact cut identity \eqref{eq:cut}, not a linearized energy balance, yields
\begin{equation}\label{eq:logistic}
 \widetilde F'-\widetilde F
 =\tau^2[a_0\widetilde F-B_\gamma\widetilde F^2]+O(\tau^3).
\end{equation}
Its positive equilibrium is $a_0/B_\gamma=F_\gamma$ from \eqref{eq:constants}.

\subsection{Uniform nonlinear trapping}\label{sec:trapping}
We detail the order of estimates, since a moving linearization without uniform trapping is insufficient. Choose a thick compact chamber with $E\in[-3,0]$, amplitude $\xi^2+\eta^2\in[F_\gamma/4,2F_\gamma]$, small fixed errors about the following five targets, and bounded auxiliary coordinates. The positivity in \eqref{eq:sign} and the graph Jacobian give a uniform analytic actual projection in this chamber, for sufficiently small $\tau$. Define
\begin{align*}
 e_s&=(C,G)-\Psi(H,b_0),& e_H&=H-H_0,\\
 e_\chi&=\chi-H\eta/b_0,& e_a&=(C_a,G_a,Q_0)-\psi_\gamma(j,\xi,\eta),\\
 e_j&=j-j_\gamma(\xi,\eta).
\end{align*}
The exact chart gives, on the thick chamber and before assuming fine errors,
\begin{align}\label{eq:hierarchy}
 e_s'&=M_se_s+O(\tau),\\
 e_H'&=e_H+\tau[f_0(H)+O(\norm{e_s})+O(\tau)],\notag\\
 e_\chi'&=O(\tau),\notag\\
 e_a'&=M_\gamma e_a+O(\norm{e_s}+|e_H|+|e_\chi|+\tau),\notag\\
 e_j'&=(1-\kappa\tau)e_j+
 O\bigl(\tau(\norm{e_s}+|e_H|+|e_\chi|+\norm{e_a})+\tau^2\bigr),\notag
\end{align}
where $f_0(H)=H^2(3b_0/2-2H)/b_0^3$. The targets have bounded $O(\tau)$ increments on the thick chamber. The leading resets of $a_d,b_d,f_a$ do not depend on their current values: those inputs enter \eqref{eq:riccati} only at orders $\tau^3,\tau^3,\tau^2$. The coarse $f$ reset is $b_0xy$.

Choose auxiliary bounds and a coarse $f$ interval first. Next choose the fine bounds
\[K_s\tau,\quad K_H\tau,\quad K_\chi\tau,\quad K_a\tau,\quad K_j\tau\]
in that order, and refine the $f$ bound last. A strict fast contraction absorbs $O(\tau)$ forcing. A slow contraction $1-c\tau$ absorbs $O(\tau^2)$ forcing. The width change $K(\tau'-\tau)=O(\tau^3)$ is smaller than every slow contraction margin. Later fine constants do not occur in earlier leading forcing constants: their input variables stay in the previously fixed thick chamber. This makes the choice noncircular.

Until a hypothetical first amplitude or charge exit, these bounds prove \eqref{eq:radialamplitudecoeff}, all physical increments are $O(t^3)$, and
\[
 |E'-E|\leq C_1\tau^3,\qquad \tau-\tau'\geq c_0\tau^3.
\]
The total charge variation is at most $(C_1/c_0)\tau_0$, by telescoping. A charge starting near $-3/2$ remains in $[-3,0]$. The leading logistic drift is uniformly inward on the protective bands $[F_\gamma/2,3F_\gamma/4]$ and $[5F_\gamma/4,3F_\gamma/2]$. Its $O(\tau^3)$ error is smaller than the order-$\tau^2$ inward drift, and a single step cannot jump across a band. The annulus and hence the actual spectral signs persist. A finite first-exit argument proves all-time trapping.

For the rational start, $0\leq\sqrt3-s<\tau_0^2$. Direct substitution gives all five initial errors $O(\tau_0^2)$, zero auxiliary coordinates, $E=-3/2+O(\tau_0)$ and $\widetilde F=1/2500+O(\tau_0)$. Uniformly for $\gamma\in[3/2,5/2]$,
\[
 F_\gamma/2<1/2500<3F_\gamma/2,
\]
with both margins greater than $10^{-6}$. Thus every sufficiently large $M$ enters strictly. Appendix~\ref{app:witness} supplies the stated explicit bound.

\subsection{Exact limit sets and angular length}
Trapping gives
\[
 t_k^{-4}=4c_0k+O(\sqrt{k}),\qquad
 \sum\tau_k^2=\infty,\quad\sum\tau_k^3<\infty.
\]
The first estimate follows by expanding reciprocal fourth powers, first obtaining equivalence and then summing the $O(t_k^2)$ remainder. Squaring $\widetilde F-F_\gamma$ in \eqref{eq:logistic} gives a stable scalar inequality with summable error and a nonsummable contraction coefficient. Hence $\widetilde F\to F_\gamma$ and $\xi^2+\eta^2\to F_\gamma$. The unwrapped phase obeys
\[
 \theta_{k+1}-\theta_k=-N\tau_k+O(\tau_k^2).
\]
Its increments eventually have one sign, tend to zero and have divergent sum. Every angle occurs as a cluster angle by the first-crossing argument.

To identify the \emph{complete} step cluster set, one needs more than $e_s=O(\tau)$. The exact chart has the form
\[
 u_{k+1}=M_su_k+g(H_k)+\tau_k r_0(E_k,\xi_k,\eta_k,\gamma)+O(\tau_k^2),
 \qquad u_k=(C_k,G_k),
\]
where $r_0$ is smooth on the fixed compact chamber. Replacing the other fast/filter variables by their targets changes this coefficient only by $O(\tau)$, hence the update by $O(\tau^2)$. The forcing changes by $O(\tau^2)$, since $\dd H=O(\tau^2)$, $\dd E=O(\tau^3)$ and $\dd(\xi,\eta)=O(\tau)$. Tracking its frozen fixed point to the next order gives $\dd C,\dd G=O(\tau^2)$. Thus
\begin{equation}\label{eq:fullstep}
 d_k=t_k^3\{c_0(E_{00}+E_{33})+c_B(E_{22}+E_{55})
 +N\eta_kF_{06}-N\xi_kF_{26}\}+O(t_k^4),
\end{equation}
where $c_B=27/256=b_0c_0$. This proves \eqref{eq:circle}; the bounded scaled coordinates give the asserted state expansion. Orthogonality of the matrix entries gives \eqref{eq:resnorm}, using
\[
 \frac{2R_\gamma^2}{\norm V_F^2/16}
 =\frac{2-\sqrt3}{\gamma(22\gamma+1)}.
\]
Radial normalization of \eqref{eq:fullstep} gives
\[
 \norm{\nu_{k+1}-\nu_k}_F
 =\frac{\sqrt3}{2}N\sqrt{F_\gamma}\,t_k^4(1+o(1)),
\]
and $t_k^4\sim(4c_0k)^{-1}$ proves \eqref{eq:resangular}. Strict entry inequalities imply a relatively open basin in the invariant block space. All other starts have the same leading state by nonexpansiveness. Block-diagonal embedding preserves the same exact dynamics in higher matrix orders. This completes Theorem~\ref{thm:resonance}.

\begin{corollary}[Residual oscillation despite a convergent speed]
For the orbit of Theorem~\ref{thm:resonance}, let $P=P_{L_\gamma}$ and $Q=I-P$, and put $H=PV$, $G=QV$. Then
\[
 \Clust\{k^{3/4}\dd y_k\}=G/4+[-R_\gamma,R_\gamma]F_{26},\qquad
 \Clust\{k^{3/4}\dd x_k\}=H/4+[-R_\gamma,R_\gamma]F_{06}.
\]
The whitened equation-residual norm has the nondegenerate cluster interval
\[
 \left[\frac{\norm G_F}{4},\
 \sqrt{\frac{\norm G_F^2}{16}+2R_\gamma^2}\right].
\]
No equality representation of this same affine space, including a redundant one, removes the scaled residual-vector oscillation.
\end{corollary}
\begin{proof}
$F_{06}\in L_\gamma$ and $F_{26}\in L_\gamma^\perp$, so applying $P,Q$ to \eqref{eq:circle} gives the two segments. The diagonal $G$ is orthogonal to $F_{26}$, whose squared Frobenius norm is two. Whitening is the exact identity in Section 3. For any representation with kernel $L_\gamma$, $\mathcal A F_{26}\ne0$. In fact $\mathcal A G$ and $\mathcal A F_{26}$ are independent, since $\mathcal A$ is injective on $L_\gamma^\perp$. The raw scaled residual has the segment $\mathcal A G/4+[-R_\gamma,R_\gamma]\mathcal A F_{26}$, and its norm is not constant on that segment.
\end{proof}

\paragraph{The nonlinear loss budget.}
For $r_k=\norm{W_{c,k}}$, the cut identity gives an exact product $r_n^2=r_K^2\exp(-L_n)$ with $L_n=\sum_{k=K}^{n-1}-\log(1-\rho_k)$ and
\[
 \rho_k=\frac{\norm{\dd W_{c,k}}^2+2v_{66,k}s_{66,k}}{r_k^2}.
\]
The proved asymptotics imply
\[
 L_n=\tfrac12\log n+O(1),\qquad
 \sum_{k=K}^{n-1}\frac{v_{66,k}s_{66,k}}{r_k^2}
 =\left(\frac14-\frac{\sqrt3}{8}\right)\log n+O(1).
\]
Thus the persistent mode is a balance of nonlinear complementarity loss and radial rescaling, not a frozen linear rotation. Theorem~\ref{thm:excess} identifies its limiting static excess as $R_\gamma^2$.

\section{The switching boundary and limits of the conclusions}
The separated sign condition fails for the earlier parameter $\gamma=1$. We give an exact frozen classification rather than extrapolating the preceding theorem.

Freeze the symmetric leading coordinates and $\chi=3\eta/4$, and let $z=(C_a,G_a,Q_0)$. Define $r,c_a,g_a$ by \eqref{eq:affinegamma} and
\[
 \Lambda=2r-Q_0-3\xi^2+\sqrt3\eta^2,\qquad
 d_*=6\gamma j+\sqrt3\eta^2.
\]
The leading actual projection across the small-eigenvalue crossing is the continuous piecewise-affine map
\begin{equation}\label{eq:switch}
 \mathcal F_\gamma(z)=
 \left(c_a+\frac32j,\
 \frac{b_0G_a+(1-b_0)g_a+(j-\xi\eta)/2}{1+b_0},\
 Q_0+3\xi^2+\max\{\Lambda,0\}-r\right).
\end{equation}
The one-step normalized error is $O(\tau)$ uniformly on compact sets, including $\Lambda=0$. Indeed the three eigenvalues of orders $t^{-1},-t^{-1},t$ have analytic separated subspaces. The remaining eigenvalue is $t^3\Lambda+O(t^5)$, and its unit eigenvector has orders $(t^2,t,t^2,1)$. Its positive part changes the leading $66$ coefficient only. Lipschitz continuity of the scalar positive part gives uniformity at zero.

Let $M_\pm$ be the two branch linear parts and $W=\diag(2,4,1)$. Direct multiplication gives, for both signs,
\[
 W-(2M_\pm-I)^TW(2M_\pm-I)
 =\frac{32(2\sqrt3-3)}{(3+8\gamma^2)^2}\ell\ell^T\succeq0,
 \quad \ell=\left(1,\frac{1-8\gamma^2}{2},-2\gamma\right)^T.
\]
The branches agree at the switching hyperplane. Splitting any connecting segment there proves global nonexpansiveness of $2\mathcal F_\gamma-I$, so the frozen map is firmly nonexpansive in this metric.

Its fixed points are parameterized by
\[
 c_f^a=2\sqrt3j+2\gamma r+\sqrt3\xi\eta,\quad
 (C_a,G_a,Q_0)=(c_f^a+3j/2,\gamma r-c_f^a/2+3j/2,r-6\gamma j),
\]
and the scalar complementarity conditions
\[
 d_*\leq0,\qquad r\geq3\xi^2,\qquad (r-3\xi^2)d_*=0.
\]
Thus there is a unique fixed point if $d_*<0$, a ray of fixed points if $d_*=0$, and no fixed point if $d_*>0$. The neutral charge has the exact increment
\begin{equation}\label{eq:neutral}
 (4\gamma C_a+Q_0)\circ\mathcal F_\gamma-(4\gamma C_a+Q_0)
 =d_*+\max\{-\Lambda,0\}.
\end{equation}
For $d_*\leq0$, every frozen orbit converges to a fixed point by the firm Fejer inequality. For $d_*>0$, every frozen orbit has drift
\[
 \frac{\mathcal F_\gamma^k(z)}k\longrightarrow
 \frac{d_*}{8\gamma^2+1}(2\gamma,0,1)^T.
\]
To prove the latter, the positive branch fixes $n=(2\gamma,0,1)^T$ linearly and contracts on its $W$-orthogonal complement. Solve its affine normal fixed equation and shift sufficiently far along $n$. This gives an exact ray orbit staying in that branch. Nonexpansiveness compares every other orbit to the ray.

At $\gamma=1$, $\xi=0$, $\eta\ne0$, the smooth proposed target $j_\gamma$ has $d_*=(\sqrt3/2)\eta^2>0$. Therefore the separated proof cannot pass through that phase: the relevant frozen fast equilibrium does not exist. The parameters move along the actual orbit, and the $O(\tau)$ error is not automatically summable. No conclusion about the originally simulated small rational start follows from this frozen classification.

The rational family nevertheless disproves the universally quantified last-step law. The universal replacement is the static-excess theorem, not a claim that every singular orbit has a limiting norm or bounded excess. The exact velocity theorem for convex forward--backward splitting uses a convex objective and a subgradient energy identity; a DRS center field need not have that structure. In the regular PSD case its leading field can be written as the saddle gradient of $\min_C\tr(H^{-1}CGC^*)$ over the affine cross-block space. The two signs, rather than convex descent of a single function, explain why a gradient-flow theorem cannot be imported without further work.

\appendix
\section{A conservative explicit rational-entry certificate}\label{app:witness}
We include quantitative bounds because ``sufficiently large rational start'' is not a specified witness. No efficiency or practical size is claimed. Let
\[
 K_{\rm aux}=10^{10000},\quad B_*=10^{100000},\quad
 C_*=B_*^{1000}=10^{100000000},\quad\tau_*=C_*^{-1000}.
\]
Use a thick chamber whose five target errors have radius at most $10^{-100}$, with $E\in[-3,0]$, amplitude in $[F_\gamma/4,2F_\gamma]$, $|f-9\sqrt3/64|\leq1$ and $|a_d|,|b_d|,|f_a|\leq K_{\rm aux}$. All other primitive parameters have modulus below $10$.

Clear $3+8\gamma^2$ from each of the six graph equations. Together with the charge equation the seven polynomials involve 24 variables. Expansion gives term counts
\[
 (14,40,40,34,38,34,34),
\]
with total degrees $(8,8,8,8,7,8,8)$ and maximum rational coefficient $16$. These finite facts follow directly from \eqref{eq:riccati} and \eqref{eq:bcharge}, and are checked by the accompanying exact symbolic script. The loose bounds of 500 terms, degree 8 and coefficients 1000 suffice below.

At $\tau=0$ the unknowns have modulus below $10$. The unscaled Jacobian determinant exceeds $1/100$ in modulus and its entries are below $100$; the adjugate bound gives inverse maximum-row-sum norm below $10^{20}$. Clearing denominators multiplies graph rows by $3+8\gamma^2\geq21$ and cannot enlarge that bound. On the complex box of coordinate modulus at most $10^{10001}$, all polynomial derivatives through order 8 are bounded by
\[
 500\cdot1000\cdot8!\,(10^{10001})^8<B_*.
\]
Set $\rho=B_*^{-4}$. On parameter polydiscs of radius $\rho$, the fixed-base-Jacobian Newton map is a contraction on the unknown ball of radius $B_*^{-2}$: its derivative is at most
\[
 24\cdot10^{20}B_*(B_*^{-2}+\rho)<1/2,
\]
and its center displacement is below $B_*^{-2}/2$. Uniform convergence gives a holomorphic graph on smaller discs. Ordinary transpose, not conjugate transpose, is used for this analytic complexification of the real matrix formulas.

For reproducible remainder bounds, set
\begin{gather*}
 D_g=1+\tau(y^2+\omega^2)+\tau^2(x^2+z^2)+\tau^3(x\omega-yz)^2,\\
 J_{00}=1+\tau\omega^2+\tau^2z^2,\quad
 J_{01}=-\tau y\omega-\tau^2xz,\quad
 J_{11}=1+\tau y^2+\tau^2x^2.
\end{gather*}
The negative spectral coefficient matrix is
$\left[\begin{smallmatrix}n_{00}/t&tn_{01}\\tn_{01}&t^3n_{11}\end{smallmatrix}\right]$, where
\begin{align*}
 n_{00}&=\frac{(B_r-\tau hy-\tau^2Kx)J_{00}
       +\tau(\eta-h\omega-\tau Kz)J_{01}}{D_g},\\
 n_{01}&=\frac{(B_r-\tau hy-\tau^2Kx)(-y\omega-\tau xz)
       +(\eta-h\omega-\tau Kz)J_{11}}{D_g},\\
 n_{11}&=\frac{(\eta+\tau(\xi x-\chi y))(-y\omega-\tau xz)
       +(-q+\xi z-\chi\omega)J_{11}}{D_g}.
\end{align*}
The physical negative part has entries
\begin{align*}
 s_{00}&=t^3(n_{00}x^2+2\tau n_{01}xz+\tau^2n_{11}z^2),\\
 s_{01}&=t^2(n_{00}xy+\tau n_{01}(x\omega+yz)+\tau^2n_{11}z\omega),\\
 s_{11}&=t(n_{00}y^2+2\tau n_{01}y\omega+\tau^2n_{11}\omega^2),\\
 s_{02}&=-t(n_{00}x+\tau n_{01}z),\qquad
 s_{12}=-(n_{00}y+\tau n_{01}\omega),\\
 s_{06}&=-t^3(n_{01}x+\tau n_{11}z),\qquad
 s_{16}=-t^2(n_{01}y+\tau n_{11}\omega),\\
 s_{22}&=n_{00}/t,\qquad s_{26}=tn_{01},\qquad s_{66}=t^3n_{11}.
\end{align*}
The second block suppresses the sixth-coordinate couplings and reverses asymmetric signs. The positive part is $R+s$.

On a smaller polydisc $|D_g-1|<1/2$, and the radial update factor $t/t'$ is within $1/2$ of one. All apparent divisions by $\tau$ in the scaled map are removable: for example $n_{00}-B_r$ and its first derivative vanish at zero, so $(n_{00}-B_r)/\tau^2$ is analytic. The largest further required divided derivative is third order, in the antisymmetric final diagonal. Cauchy estimates on successively halved discs bound the scaled updates by $B_*^{30}$, the charge update by $B_*^{70}$, and compositions with the targets and quartic correction by $B_*^{100}$. Derivatives and Lipschitz constants through the required orders are below $B_*^{200}<C_*$. The gaps in these exponents absorb all finite dimension factors and equivalence of the displayed weighted norms.

Leading auxiliary resets are independent of their auxiliary inputs. Applying the same estimate with zero such inputs and $B_{\rm small}=10^{30}$ bounds them by $10^{6000}<K_{\rm aux}/2$. Hence their coarse box is invariant for sufficiently small $\tau$.

Use the fine bounds
\[
 \norm{e_s}\leq C_*^2\tau,\quad |e_H|\leq C_*^4\tau,\quad
 |e_\chi|\leq C_*^6\tau,\quad \norm{e_a}\leq C_*^{10}\tau,
 \quad |e_j|\leq C_*^{14}\tau,
\]
and $|f-9\sqrt3/64|\leq C_*^8\tau$. At $\tau\leq\tau_*$ these lie well inside the thick chamber. The corresponding forcing bounds, in the order of \eqref{eq:hierarchy}, are
\[
 C_*\tau,\quad 2C_*^3\tau^2,\quad C_*\tau,\quad
 4C_*^7\tau,\quad5C_*^{11}\tau^2,
\]
whereas the fast margins are $10^{-4}$ times their fine radii and the slow margins are at least $\tau$ times their radii. The refined $f$ forcing is at most $3C_*^5\tau$. Every forcing and width change is below half the corresponding margin.

The physical step bound is $\norm d_F\leq C_*^{20}t^3$, so $|\dd E|\leq C_*^{42}\tau^3$. Its total variation is below $10C_*^{42}\tau_0<1/10$. The corrected-amplitude remainder is at most $C_*^{50}\tau^3$, while the inward drift on the protective bands exceeds $10^{-7}\tau^2$. The small-eigenvalue margin exceeds $10^{-6}$ and the initial amplitude gap exceeds $10^{-6}$. All perturbations are smaller than half these margins for $\tau\leq C_*^{-1000}$.

Finally, direct substitution of the rational start bounds its five errors and refined $f$ error by $10^6\tau_0^2$, and the charge and amplitude errors by $10^6\tau_0$. Thus
\[
 M\geq C_*^{500}=10^{50000000000}
 \quad\Longrightarrow\quad \tau_0=M^{-2}\leq\tau_*.
\]
Every entry and trapping inequality follows. In particular $M=10^{10^{12}}$ specifies one exact rational witness for every rational parameter in the displayed interval. This is an analytic existence certificate, not a finite simulation or a practical complexity estimate.

\section{Reproducibility and audit boundary}
Three accompanying scripts use exact SymPy arithmetic:
\begin{center}\texttt{gamma\_family\_algebra.py}\\
\texttt{gamma\_spectral\_chart.py}\\
\texttt{switching\_fast\_map.py}.\end{center} They verify the affine fast fixed point, Gram-defect minors, filter identity, quartic correction, Schur determinant, reconstruction of the weighted Riccati equations from the original reflected matrix, low-order spectral coefficients, and the switching neutral-charge identities. They do not replace the uniform analytic and first-exit arguments.

The main mathematical dependencies are: semialgebraic one-variable Puiseux asymptotics; elementary nonexpansive comparison; cone projection complementarity; analytic matrix spectral subspaces separated by sign; and the explicit estimates displayed above. Historical priority remains a separate question. A fixed publication sample has been selected from two 2024--2025 issues per journal, but its full-proof comparisons and any defensible upper-quintile assessment are not complete. No probability of editorial acceptance is claimed.

\begin{thebibliography}{9}
\bibitem{BGV} J. Bolte, S. Gaubert and G. Vigeral, \emph{Definable zero-sum stochastic games}, Mathematics of Operations Research \textbf{40} (2015), 171--191. DOI: 10.1287/moor.2014.0666. Primary preprint: arXiv:1301.1967v2, Theorem 3.
\bibitem{LRY} Y. Liu, E. K. Ryu and W. Yin, \emph{A new use of Douglas--Rachford splitting and ADMM for identifying infeasible, unbounded, and pathological conic programs}, Mathematical Programming \textbf{177} (2019), 225--253. Primary preprint: arXiv:1706.02374v3, Theorems 2.3, 2.5 and 2.6.
\bibitem{DJK} T. Do, A. Jentzen and B. Kuckuck, \emph{On the gradient conjecture at infinity: Finite spherical length for gradient trajectories of definable potentials}, arXiv:2610.29211v2, revised 9 November 2026, Theorem 1.1.
\bibitem{Ryu} E. K. Ryu, \emph{Cosmic divergence, weak cosmic convergence, and fixed points at infinity}, arXiv:1709.02838. This concerns the unrestricted cosmic-convergence question, not the rational affine--PSD last-step counterexample proved here.
\end{thebibliography}
\end{document}
